% Source: https://artofproblemsolving.com/community/c3h1385802_not_as_easy_as_it_looks
% Stable identity: aops-topic-1385802-q01
Let $x$ be the area of $DBE$, so triangle $CDB$ has area $40+x$.

\begin{asy}
usepackage("amsmath");
size(6.5cm);
defaultpen(fontsize(10pt));
pair ft(pair P,pair U,pair V){return U+dot(P-U,V-U)/dot(V-U,V-U)*(V-U);}
path ra(pair P,pair u,pair v,real s=0.2){return P+s*unit(u)--P+s*unit(u)+s*unit(v)--P+s*unit(v);}
real a=1, bb=sqrt(5);
pair C=(0,0),B=(bb,0),A=(0,a);
pair D=ft(C,A,B); pair Ee=(D.x,0);
draw(A--B--C--cycle,linewidth(0.9));
draw(C--D); draw(D--Ee);
draw(ra(C,B-C,A-C,0.07)); draw(ra(D,B-D,C-D,0.07)); draw(ra(Ee,B-Ee,D-Ee,0.07));
label("$A$",A,N); label("$B$",B,SE); label("$C$",C,SW); label("$D$",D,NE); label("$E$",Ee,S);
label("$48$",(A+C+D)/3+(-0.03,0.02),fontsize(9pt));
label("$40$",(C+D+Ee)/3+(0.03,0),fontsize(9pt));
label("$x$",(D+Ee+B)/3+(0.05,0),fontsize(9pt));
label("\mbox{}",(min(currentpicture,user=true)+max(currentpicture,user=true))/2,black); // keeps the text after the figure black when LaTeX typesets the labels inline
\end{asy}

Triangles $ACD$ and $CBD$ are similar (both are right-angled at $D$ and $\angle ACD=90^\circ-\angle A=\angle CBD$), with $CD$ corresponding to $BD$ and $AD$ to $CD$. Similarly, in right triangle $CDB$ the altitude $DE$ creates similar triangles $CDE$ and $DBE$, with $CD$ corresponding to $DB$. Areas of similar triangles are in the ratio of the squares of corresponding sides, so
\[
\dfrac{[ACD]}{[CBD]}=\left(\dfrac{CD}{BD}\right)^2=\dfrac{[CDE]}{[DBE]} .
\]
Hence $\dfrac{48}{40+x}=\dfrac{40}{x}$, so $48x=1600+40x$ and $x=\boxed{200}$.
